\documentclass[11pt]{article}

\usepackage[a4paper,margin=1in]{geometry}
\usepackage{amsmath,amssymb}
\usepackage{hyperref}

\hypersetup{
    colorlinks=true,
    linkcolor=blue,
    citecolor=blue,
    urlcolor=blue
}

\title{Why $W_2(\delta_0,\delta_1)=1$}
\author{}
\date{}

\begin{document}

\maketitle

Let
\[
\mu=\delta_0,
\qquad
\nu=\delta_1.
\]
By definition,
\[
W_2(\mu,\nu)
=
\left(
\inf_{\gamma\in\Pi(\mu,\nu)}
\int |x-y|^2\,d\gamma(x,y)
\right)^{1/2}.
\]

Since all mass starts at $0$ and must end at $1$, the only possible transport
plan is
\[
\gamma=\delta_{(0,1)}.
\]
Thus
\[
\int |x-y|^2\,d\gamma(x,y)
=
|0-1|^2
=
1.
\]
Therefore,
\[
W_2^2(\delta_0,\delta_1)=1,
\qquad
W_2(\delta_0,\delta_1)=\sqrt{1}=1.
\]

More generally,
\[
W_p(\delta_a,\delta_b)=|a-b|.
\]

\end{document}
